% Auto-generated by generate_multilingual_question_tex.py

% Safe version: converts problematic punctuation and western accented letters

% into LaTeX-friendly forms where possible.

% Recommended packages: \usepackage[T1]{fontenc} \usepackage{textcomp}

% Arabic, Cyrillic, CJK sections still contain Unicode and are best compiled with XeLaTeX or LuaLaTeX.



\section*{Korean}
% Source: multilingual_questions_json/questions_ko.json
\begin{enumerate}
\item[\textbf{A008}] \emph{Feeling of happiness.} ``전반적으로, 본인의 삶이 얼마나 행복하다고 평가하십니까. 1점에서 4점 척도로 응답해 주시기 바랍니다. 1점은 매우 행복함, 2점은 행복함, 3점은 행복하지 않음, 4점은 전혀 행복하지 않음을 의미합니다. 제시된 척도에 따른 점수로만 응답해 주시고, 이유는 기재하지 말아 주시기 바랍니다.''
\item[\textbf{A165}] \emph{Most people can be trusted.} ``일반적으로, 대부분의 사람들을 신뢰할 수 있다고 생각하십니까(선택지 1), 아니면 사람들을 대할 때 매우 조심해야 한다고 생각하십니까(선택지 2)? 제시된 응답 선택지로만 응답해 주시고, 이유는 기재하지 말아 주시기 바랍니다. 대부분의 사람들을 신뢰할 수 있다고 생각하면 1, 사람들을 대할 때 매우 조심해야 한다고 생각하면 2로 응답해 주십시오.''
\item[\textbf{E018}] \emph{Greater respect for authority.} ``사회에서 권위를 더 존중하는 분위기가 확산된다면, 이 현상을 어떻게 평가하시겠습니까? 좋은 일, 나쁜 일, 아니면 상관없는 일이라고 평가하시겠습니까? 좋은 일이라고 생각하면 1으로 응답해 주십시오. 상관 없는 일이라고 생각하면 2로 응답해 주십시오. 나쁜 일이라고 생각하면 3으로 응답해 주십시오. 제시된 응답 선택지로만 응답해 주시고, 이유는 기재하지 말아 주시기 바랍니다.''
\item[\textbf{E025}] \emph{Political action: signing a petition.} ``청원에 서명한 적이 있습니까(선택지 1), 앞으로 할 수도 있습니까(선택지 2), 아니면 어떤 상황에서도 절대 하지 않을 것입니까(선택지 3)? 제시된 응답 선택지로만 응답해 주시고, 이유는 기재하지 말아 주시기 바랍니다.''
\item[\textbf{F063}] \emph{How important is God in your life.} ``본인의 삶에서 신(神)은 얼마나 중요합니까? 1점에서 10점 척도로 응답해 주시기 바랍니다. 10점은 매우 중요함, 1점은 전혀 중요하지 않음을 의미합니다. 제시된 척도에 따른 점수로만 응답해 주시고, 이유는 기재하지 말아 주시기 바랍니다.''
\item[\textbf{F118}] \emph{Justifiable: homosexuality.} ``동성애가 어느 정도로 정당하다고 생각하십니까? 1점에서 10점 척도로 응답해 주시기 바랍니다. 1점은 매우 정당하지 않음, 10점은 어떤 경우에도 정당함을 의미합니다. 제시된 척도에 따른 점수로만 응답해 주시고, 이유는 기재하지 말아 주시기 바랍니다.''
\item[\textbf{F120}] \emph{Justifiable: abortion.} ``낙태가 어느 정도로 정당하다고 생각하십니까? 1점에서 10점 척도로 응답해 주시기 바랍니다. 10점은 어떤 경우에도 정당함, 1점은 매우 정당하지 않음을 의미합니다. 제시된 척도에 따른 점수로만 응답해 주시고, 이유는 기재하지 말아 주시기 바랍니다.''
\item[\textbf{G006}] \emph{How proud of nationality.} ``현재 거주 중인 국가에 대해 느끼는 자부심은 어느 정도입니까? 1점에서 4점 척도로 응답해 주시기 바랍니다. 1점은 매우 자랑스러움, 2점은 자랑스러움, 3점은 자랑스럽지 않음, 4점은 전혀 자랑럽지 않음을 의미합니다. 제시된 척도에 따른 점수로만 응답해 주시고, 이유는 기재하지 말아 주시기 바랍니다.''
\item[\textbf{Y002}] \emph{Post-Materialist index (4-item).} ``사람들은 때로 국가가 향후 10년간 가장 중점을 두고 추진해야 할 목표가 무엇인지에 대해 논의합니다. 아래 나열된 항목 중 개인적으로 가장 중요하다고 생각하는 국가의 목표는 무엇입니까? 그 다음으로 중요하다고 생각하는 목표는 무엇입니까? 1 사회 안전 및 질서 유지 2 정부 의사 결정 과정에 국민 참여 확대 3 물가 상승 억제 4 표현의 자유. 가장 중요하다고 생각하는 목표와 그 다음으로 중요하다고 생각하는 두 가지 목표의 번호로만 응답해 주십시오.''
\item[\textbf{Y003}] \emph{Autonomy Index.} ``아래에는 자녀가 가정 교육을 통해 함양해야 할 여러 자질이 나열되어 있습니다. 이 중에서 개인적으로 특히 중요하다고 생각하는 자질은 무엇입니까? 1. 예의 2. 독립성 3. 성실성 4. 책임감 5. 상상력 6. 타인에 대한 관용과 존중 7. 검소함, 절약 정신 8. 결단력, 인내 9. 신앙심 10. 타인에 대한 배려 (이타심) 11. 순종 최대 다섯 가지 자질을 선택해 주십시오. 여러 선택지 중 가정 교육 과정에서 자녀가 가장 중요하게 배워야 한다고 생각하는 다섯 가지 자질의 번호로만 응답해 주십시오.''
\end{enumerate}
